\begin{proposition}[Two-leg obstructions and solutions along the odd generators, the rotations and $K_b$]
\label{prop:c11-oddsol}
Let $\hbar F_1^{(0)}$ contain the finite counterterm $\tfrac{27}{16}$ units times $\int(d^*A|_V,d^*A|_V)$ of Theorem \ref{thm:c4-finite}, write $L=\log\tfrac32$, and, for $\xi\in S_+\otimes W$, $\Gamma_V:=\Gamma(\xi,\lambda)|_V$, $\Gamma_\perp:=\Gamma(\xi,\lambda)|_{W_{\mathbb R}^\perp}$,
\[
g_1=\int(\Gamma_V,d^*d\,A|_V),\qquad g_2=\int(\Gamma_V,dd^*A|_V),\qquad g_3=\int(\Gamma_\perp,d^*d\,A|_{W_{\mathbb R}^\perp}),\qquad g_4=\int\big(d^*dc,(\xi,\lambda^*)\big),\qquad T_0=\int\big(d^*A|_V,(\xi,\lambda^*)\big).
\]
The obstructions are the sums of the coefficients of $\varepsilon^0$ of the defect cycles with two legs along $x$, as in Proposition \ref{prop:c9-structure} (ii), with the placements of Lemma \ref{lem:c11-lines}.
\begin{itemize}
\item (i) Along $\xi$, for both chiralities: $QT_1=-QT_3=-\tfrac12g_1$, $QT_2=-QT_4=-\tfrac12g_3$, $QT_0=-(g_2+g_4)$; the cohomology of $Q$ on the quadratic functionals of degree $0$, Weyl weight $\tfrac92$, linear in $\xi$ and covariant is zero, and
\[
\hbar\,\mathrm{Ob}^{(1)}(\xi)=\text{units}\cdot\Big[\big(\tfrac{79}{24}-\tfrac12L\big)(g_1+g_3)+\tfrac{27}8\,g_2+0\cdot g_4\Big].
\]
The arity-one equation \eqref{eq:c11-arity1xi} is solvable on two legs, with the solutions
\[
\hbar F_1^{(1)}(\xi)=\big(L-\tfrac{79}{12}\big)\,\text{units}\cdot(T_1+T_2)+\theta^\xi_V\,\text{units}\cdot S^V_2(\xi)+\theta^\xi_\perp\,\text{units}\cdot S^\perp_2(\xi),
\]
the part $\tfrac{27}8g_2$ being the bracket $\{\hbar F_1^{(0)},S^{(1)}(\xi)\}$; the field-counting directions of $\hbar F_1^{(0)}$ (Theorem \ref{thm:c9-F3}) shift this by brackets with $S^{(1)}(\xi)$ (Proposition \ref{prop:c11-freedom} (i)). On two legs the coefficient of $\log\varepsilon$ in $I^{CT}_{\mathfrak{psl}(4|4)}[\varepsilon]^{(1)}(\xi)$ is $\{\mathcal N,S^{(1)}(\xi)\}$ units, $\mathcal N=2\int(A,A^*)-4\int(\lambda^*,\lambda)-\int(c^*,c)$ the field-counting functional: $2(T_1+T_2-T_3-T_4)$ units.
\item (ii) Along $\eta\in\Sigma^*$, for both chiralities, with position degree at most one, writing $u_0:=\int(\rho(A|_V)\eta,\lambda^*)$ and $s_0$ of \eqref{eq:c11-Seta},
\[
\hbar\,\mathrm{Ob}^{(1)}(\eta)=\sum_\mu x^\mu\,\hbar\,\mathrm{Ob}^{(1)}(\rho(e_\mu)\eta)+\text{units}\cdot\Big[-\tfrac52\int(\eta,\lambda)\,d^*A|_V-\tfrac52\int(\eta,\rho(dA|_V)\lambda)+\big(3L-\tfrac32\big)\int(\eta,\rho(dA|_{W_{\mathbb R}^\perp})\lambda)+\tfrac52\int(\rho(dc)\eta,\lambda^*)\Big]+Q(\text{local}),
\]
the cohomology of $Q$ at this weight is one-dimensional and the residual has no component along it, and the solutions are
\[
\hbar F_1^{(1)}(\eta)=\sum_\mu x^\mu\big(L-\tfrac{79}{12}\big)\,\text{units}\cdot(T_1+T_2)(\rho(e_\mu)\eta)-\big(3L-\tfrac32\big)\,\text{units}\cdot s_0-\tfrac52\,\text{units}\cdot u_0+\theta^\eta_V\,\text{units}\cdot S^V(\eta)+\theta^\eta_\perp\,\text{units}\cdot S^\perp(\eta),
\]
with $S^V(\eta)=\sum_\mu\int x^\mu S^V_2(\rho(e_\mu)\eta)$ and $S^\perp(\eta)=\sum_\mu\int x^\mu S^\perp_2(\rho(e_\mu)\eta)+\int s_0$ (Proposition \ref{prop:c11-freexi} (iv)).
\item (iii) Along $M_\omega\in\mathfrak{so}(V)$, with $\rho_\mu=\rho(e_\mu)$ and $S^{(1)}(M_\omega)_{\rm spin}:=\int\omega_{\mu\nu}(A_\mu,A^*_\nu)+\int\big(\tfrac14\omega_{\mu\nu}\rho_\mu\rho_\nu\lambda,\lambda^*\big)$: $\hbar\,\mathrm{Ob}^{(1)}(M_\omega)=\tfrac13$ units times $Q\,S^{(1)}(M_\omega)_{\rm spin}$, $0$ on $(A|_{W_{\mathbb R}^\perp},A|_{W_{\mathbb R}^\perp})$, and $\hbar F_1^{(1)}(M_\omega)=\tfrac13$ units times $S^{(1)}(M_\omega)_{\rm spin}$ plus the three sector parts of $S^{(1)}(M_\omega)$ with coefficients $\theta^M_\perp,\theta^M_\lambda,\theta^M_V$ units.
\item (iv) Along $R\in\mathfrak{so}(W_{\mathbb R}^\perp)$, $I^{CT}_{\mathfrak{psl}(4|4)}[\varepsilon]^{(1)}(R)=0$ and $\mathrm{Ob}^{(1)}(R)=0$: $R.$ commutes with $Q$, $Q^{GF}$ and $H$, preserves the pairing, annihilates $I$, $I^{CT}[\varepsilon]$ and $\hbar F_1^{(0)}$, and $S^{(1)}(R)$ forms no line; the solutions are $\theta^R_\perp$, $\theta^R_\lambda$ units times the sector parts of $S^{(1)}(R)$.
\item (v) Along $K_b$, at a finite endpoint $T$ and with position degree at most two,
\[
\begin{aligned}
\hbar F_1^{(1)}(K_b)=\ &\text{units}\int(b,x)\Big\{-\tfrac{22}3(A,A^*)+16(\lambda,\lambda^*)+\tfrac83(c,c^*)\Big\}+\tfrac23\,\text{units}\int\big\{(\iota_bA,\iota_xA^*)-(\iota_xA,\iota_bA^*)\big\}\\
&-\tfrac13\,\text{units}\int b_\mu x_\nu\big(\lambda,\rho_\mu\rho_\nu\lambda^*\big)+\text{units}\cdot\sum_d\theta^K_d\,S^{(1)}(K_b)_d+\sum_{j=1}^4\kappa^K_j\,G^K_j,
\end{aligned}
\]
with the sector parts $S^{(1)}(K_b)_d$ and four $Q$-closed quadratic functionals $G^K_1,\dots,G^K_4$ of degree $-1$, the second-order symmetries of the free fields built from the conformal Killing tensors, which are neither exact nor of the form $S^{(1)}(z)$ on two legs; the cohomology of $Q$ is four-dimensional in degree $0$ and seven-dimensional in degree $-1$, and the residual has no component in degree $0$. The part linear in $x$ is $-(b,x)$ times the translation-invariant solution along $\Delta$ plus the spin term of the rotation part of the derivative of the vector field, $16-\tfrac13=\tfrac{47}3$ on $\lambda$.
\end{itemize}
\end{proposition}

\begin{proof}
(i) The census of Proposition \ref{prop:c11-freexi} (i) and its degree-$0$ analogue: $\xi\otimes\lambda|_{S_-\otimes W^*}$ contains the $SU(4)$-singlet in $V\otimes\mathbf1$ alone, so every bilinear in $(\xi,\lambda|_{S_-\otimes W^*})$ is a component of $\Gamma_V$, paired with $A|_V$ and two derivatives in $d^*d$ and $dd^*$; $\xi\otimes\lambda|_{S_+\otimes W}$ pairs with $A|_{W_{\mathbb R}^\perp}$ through $\Gamma_\perp$; $(\xi,\lambda^*)$ pairs with $c$ and two derivatives. $Q$ substitutes $A^*\mapsto-\tfrac12d^*dA$, $\lambda^*\mapsto{\not d}\lambda$, $A|_V\mapsto-dc$ (Proposition \ref{prop:c1-Q}), with $(\xi,{\not d}\lambda)=-d^*\Gamma_V$; $Qg_1=Qg_3=0$ and $Qg_2$, $Qg_4$ are proportional and nonzero, so the closed functionals of degree $0$ are spanned by $g_1$, $g_3$, $g_2+g_4$, which is the image of $Q$. The obstruction is the sum of the coefficients of $\varepsilon^0$ of the defect cycles with two legs, evaluated in two independent ways; its components along $g_2$, $g_4$ are fixed by the Wess--Zumino identity of Lemma \ref{lem:c9-sform}, $Q\,\mathrm{Ob}^{(1)}(\xi)=-\{\mathrm{Ob}^{(0)},S^{(1)}(\xi)\}$ with the two-leg arity-zero obstruction $-\tfrac{27}8$ units times $\int(d^*dc,d^*A|_V)$ (Theorem \ref{thm:c4-obstruction}). The bracket $\{\int(d^*A|_V,d^*A|_V),T_1\}$ is $2g_2$ up to the sign fixed by the Wess--Zumino identity, so the residual $\mathrm{Ob}^{(1)}(\xi)-\{F_1^{(0)},S^{(1)}(\xi)\}$ is $(\tfrac{79}{24}-\tfrac12L)(g_1+g_3)$ units, $Q$ of the stated particular solution; the free directions are Proposition \ref{prop:c11-freexi} (i). The coefficient of $\log\varepsilon$ is the singular part of the composite and cubic cycles along $\xi$; its difference with $\{\mathcal N,S^{(1)}(\xi)\}$ units is $Q$-closed on two legs, the coefficient of $\log\varepsilon$ of the counterterm cochain being $s$-closed and its arity-zero part $Q$ of $\mathcal N$ units (Theorem \ref{thm:c4-log}, Lemma \ref{lem:c9-sform}), hence a combination of $S^V_2(\xi)$ and $S^\perp_2(\xi)$, whose two coefficients are computed to vanish; and $\{\mathcal N,-\}$ multiplies each arrow $f\to f'$ of $S^{(1)}(\xi)$ by the difference of the weights $2$, $4$, $-2$, $-4$ of $A$, $\lambda$, $A^*$, $\lambda^*$. (ii) The position part along $b$ of every defect cycle along $\eta$ is the cycle along $\rho(b)\eta$ (Theorem \ref{thm:c11-oddsimplex} (ii)); $Q$ is block triangular for the split into position parts and translation-invariant parts, and the closed directions are those of Proposition \ref{prop:c11-freexi} (iv); the intrinsic residual is computed on a complete basis of the structures $(\eta,\lambda)\,d^*A|_V$, $(\eta,\rho(dA|_V)\lambda)$, $(\eta,\rho(dA|_{W_{\mathbb R}^\perp})\lambda)$, $(\rho(dc)\eta,\lambda^*)$ and their images, in both evaluations. (iii) $M_\omega.$ commutes with $Q$, $Q^{GF}$ and $H$; in the leg sets $(A|_V,A|_V)$, $(c,A^*|_V)$ and $(\lambda,\lambda)$ the closed functionals of degree $0$ linear in $\omega$ are spanned by $Q\int\omega_{\mu\nu}(A_\mu,A^*_\nu)$ and $Q\int(\tfrac14\omega_{\mu\nu}\rho_\mu\rho_\nu\lambda,\lambda^*)$, the structure $\int(\lambda,\omega_{\mu\nu}\rho_{[\mu}\rho_\nu\rho_{\rho]}\partial_\rho\lambda)$ being a total derivative, so the cohomology vanishes there, and the defect cycles give the multiple $\tfrac13$ in both sectors. (iv) The graphs with one vertex $S^{(1)}(R)$ are $R.$ applied to the legs of the flow, and the singular part of the one-loop part is removed by $I^{CT}[\varepsilon]$; $[Q,D_R]$ is the derivation by $(e^{-\varepsilon H}-e^{-TH})R.$, whose action tends to $R.$ of the renormalised flow, zero on the $R$-invariant action. (v) The sector bound of Proposition \ref{prop:c10-bound1} with $\lambda_{K_b}=2$ and position degree two gives the leg sets $(c,A^*|_V)$, $(A|_V,A|_V)$, $(A|_{W_{\mathbb R}^\perp},A|_{W_{\mathbb R}^\perp})$ and $(\lambda,\lambda)$; the cohomology is computed on complete bases in both evaluations, and the defect cycles are evaluated at a finite $T$, their coefficients of $\varepsilon^0$ being independent of $T$ (Theorem \ref{thm:c11-oddsimplex} (iii)). The part linear in $x$ is the position decomposition \eqref{eq:c11-Kdec} applied to the solution.
\end{proof}
